% No modificar estas líneas de código, por favor dirigirse a MARCO PRÁCTICO 

\newpage

%------------------------
%
%       MARCO PRÁCTICO
%
%------------------------

\section{MARCO PRÁCTICO}
Introducir lo que se desea.... ahora aprovecharemos para demostrar como adicionar notas a pie de página, siendo estas útiles para ciertas circunstancias a ser determinadas por los autores.

Aunque dicen que no es conveniente poner demasiadas, para no sobrecargar el texto, también se pueden poner notas\footnote{Una nota a pie de página.} a pie de página. 

Para ello escribimos \verb|\footnote{texto de la nota}|. 

Aquí va otra nota\footnote{Otra nota más.} a pie de página solo para demostrar lo explicado anteriormente y su funcionalidad. 

Para formalidades, cuando se desea utilizar comillas dobles, no olvidemos que debemos abrirlas y luego cerrarlas, en algunos casos el \LaTeX no reconoce las mismas, es por eso que se debe utilizar de la siguiente forma: \verb|``click''|, siendo el resultado a obtener el siguiente: ``click''. 